% =============================================================================
% CESE5040 Lab report -- generic template.
%
% TO INSTANTIATE FOR A NEW LAB:
%   1. Copy this whole template/ folder to ../labN/  (use ../new-lab.sh N).
%   2. Rename Lab.tex -> LabN.tex.
%   3. Edit the THREE \renewcommand lines below to set lab number and (if a
%      different student is using the template) name + number.
%   4. Edit \labtitle to match the lab's actual subtitle.
%   5. Add or remove "Question N" subsections to match the lab's exercises.
%
% Compile: pdflatex LabN.tex   (twice, for cross-references)
% Submit:  rename produced PDF to LabN_Firstname_Lastname_StudentNumber.pdf
% =============================================================================

\documentclass[11pt, a4paper]{article}

% ---------------------------------------------------------------------------
% Lab-specific identity. Edit these for each lab.
% ---------------------------------------------------------------------------
\newcommand{\labtitle}{Lab Assignment N: <subtitle>}

% =============================================================================
% Shared preamble for CESE5040 lab reports.
% Lab-agnostic. Copy this file into each labN/ folder so each lab is a
% self-contained Overleaf-uploadable project.
%
% Compile with pdflatex.
% =============================================================================

% ---------- Page geometry ----------------------------------------------------
\usepackage[a4paper, margin=2.5cm, headheight=14pt]{geometry}

% ---------- Fonts and language -----------------------------------------------
\usepackage[T1]{fontenc}
\usepackage[utf8]{inputenc}
\usepackage{lmodern}
\usepackage[english]{babel}
\usepackage{microtype}

% ---------- Math -------------------------------------------------------------
\usepackage{amsmath, amssymb, amsthm}
\usepackage{mathtools}
\usepackage{bm}
\usepackage{siunitx}
\sisetup{
    detect-weight        = true,
    detect-family        = true,
    range-phrase         = {--},
    range-units          = single,
    per-mode             = symbol,
}

% ---------- Tables and figures -----------------------------------------------
\usepackage{graphicx}
\usepackage{booktabs}
\usepackage{array}
\usepackage{multirow}
\usepackage{caption}
\captionsetup{labelfont=bf, font=small, skip=6pt}
\usepackage{subcaption}
\usepackage{float}

% ---------- Code listings ----------------------------------------------------
\usepackage{listings}
\usepackage{xcolor}
\definecolor{pyKeyword}{HTML}{0033B3}
\definecolor{pyString}{HTML}{067D17}
\definecolor{pyComment}{HTML}{8C8C8C}
\definecolor{pyNumber}{HTML}{1750EB}
\definecolor{pyBg}{HTML}{F7F7F7}
\definecolor{pyFrame}{HTML}{D0D0D0}

\lstdefinestyle{python}{
    language            = Python,
    basicstyle          = \ttfamily\footnotesize,
    keywordstyle        = \color{pyKeyword}\bfseries,
    stringstyle         = \color{pyString},
    commentstyle        = \color{pyComment}\itshape,
    numberstyle         = \tiny\color{pyComment},
    numbers             = left,
    numbersep           = 8pt,
    stepnumber          = 1,
    breaklines          = true,
    breakatwhitespace   = true,
    showstringspaces    = false,
    frame               = single,
    rulecolor           = \color{pyFrame},
    backgroundcolor     = \color{pyBg},
    captionpos          = b,
    tabsize             = 4,
    columns             = fullflexible,
    keepspaces          = true,
    morekeywords        = {self, np, plt, jit, prange},
}

\lstdefinestyle{shell}{
    basicstyle          = \ttfamily\footnotesize,
    breaklines          = true,
    showstringspaces    = false,
    frame               = single,
    rulecolor           = \color{pyFrame},
    backgroundcolor     = \color{pyBg},
    captionpos          = b,
}

\lstset{style=python}

% ---------- Hyperlinks and cross-refs ----------------------------------------
\usepackage[hidelinks]{hyperref}
\usepackage[capitalize, noabbrev]{cleveref}

% ---------- Headers and footers ----------------------------------------------
\usepackage{fancyhdr}
\pagestyle{fancy}
\fancyhf{}
\fancyhead[L]{\small CESE5040 -- HPC and AI Architectures}
\fancyhead[R]{\small \labnumber{} -- \studentname{} (\studentnumber{})}
\fancyfoot[C]{\thepage}
\renewcommand{\headrulewidth}{0.4pt}
\renewcommand{\footrulewidth}{0pt}

% ---------- Section formatting -----------------------------------------------
\usepackage{titlesec}
\titleformat{\section}{\Large\bfseries}{\thesection}{0.7em}{}
\titleformat{\subsection}{\large\bfseries}{\thesubsection}{0.6em}{}
\titleformat{\subsubsection}{\normalsize\bfseries}{\thesubsubsection}{0.5em}{}

% ---------- Identity placeholders --------------------------------------------
% Override these in the main file via \renewcommand BEFORE \begin{document}.
\newcommand{\labnumber}{Lab~?}
\newcommand{\studentname}{Firstname Lastname}
\newcommand{\studentnumber}{0000000}

% ---------- Authoring helpers ------------------------------------------------
% Question stub matching the docx template.
\newcommand{\answerhere}{\textit{Give answer here\dots}}
\newcommand{\todoitem}[1]{\textcolor{red}{\textbf{TODO:} #1}}

% Inline code shorthand
\newcommand{\code}[1]{\texttt{#1}}

% Big-O notation
\newcommand{\bigO}[1]{\ensuremath{\mathcal{O}\!\left(#1\right)}}


\renewcommand{\labnumber}{Lab~N}
\renewcommand{\studentname}{Daniel Tyukov}
\renewcommand{\studentnumber}{5714699}

% ---------------------------------------------------------------------------

\title{%
    \vspace{-2em}
    \large CESE5040 -- High-Performance Computing and AI Architectures \\[0.4em]
    \Large \textbf{\labtitle} \\[0.6em]
    \rule{\linewidth}{0.4pt}
}
\author{\studentname{} \\ \texttt{\studentnumber}}
\date{\today}

\begin{document}
\maketitle
\thispagestyle{fancy}

% =============================================================================
\section{Student information}
% =============================================================================

\noindent\textbf{Name:}~\studentname{}\\
\textbf{Student number:}~\studentnumber{}

\paragraph{Notes.}
\begin{itemize}
    \itemsep0em
    \item Provide your answers following the template made available in Brightspace.
    \item Provide your name and student number inside the report.
    \item Provide your student number on the file name.
    \item Provide your answers at least 3 days before the next lecture.
    \item When answering the questions, take all measurements using the same machine.
\end{itemize}

% =============================================================================
\section{Lab Exercise N}
% =============================================================================

% Add or remove Question subsections to match the actual lab's exercise count.
% Inside each, use \paragraph{(a) ...} for sub-questions when needed.

\subsection{Question 1}
\answerhere{}

\subsection{Question 2}
\answerhere{}

\subsection{Question 3}
\answerhere{}

% Example reusable building blocks ------------------------------------------

% Timing table:
% \begin{table}[H]
%     \centering
%     \caption{Execution times for <thing> over <K> timesteps on the course AWS VM.}
%     \label{tab:timings}
%     \begin{tabular}{lrr}
%         \toprule
%         Variant & Time (s) & Speedup ($\times$) \\
%         \midrule
%         Baseline      &  & 1.00 \\
%         Optimized     &  &      \\
%         \bottomrule
%     \end{tabular}
% \end{table}

% Code listing from your submitted .py file:
% \lstinputlisting[style=python,
%                  caption={\code{LN\_Qx.py} -- <description>.},
%                  label={lst:qx}]{code/LN_Qx.py}

% Figure:
% \begin{figure}[H]
%     \centering
%     \includegraphics[width=0.7\linewidth]{figures/<plot>.pdf}
%     \caption{<caption>.}
%     \label{fig:<plot>}
% \end{figure}

% =============================================================================
\section{LLM Usage Disclosure}
% =============================================================================

\noindent
Disclose your usage of LLM for this assignment. Refer to the lecture slides and
reporting guidelines to find out the LLM usage policy of the course.

\paragraph{Example.}
\textit{I used ChatGPT to help with NumPy documentation and refine the language of my report.}

\paragraph{My disclosure.}
\answerhere{}

\end{document}
